\section{Transporting the equation ideal through all spectator tests}
\label{sec:actual-equation-kernel-relay}

The finite measurements used for the rank bound must be supplemented by
an assertion about every spectator word. We prove that assertion on the
selected history. The mechanism has two parts: the marked section obeys
an exact microscopic equation identity, and its comparison with the
physical section has a vanishing centered discrepancy. The same
calibration map occurs in both parts.

We retain the funded calibrated family for a simple gauge group and its
centered observable bounds. For each fixed grade, assume the selected
local reserve, terminal coefficient, response and finite-census purchases
hold eventually in the mesh cutoff, uniformly for every period, on this
same family. The derivative credit and sphere bounds are the purchases
of Section~\ref{sec:ex-marked-portfolio-calibration}. The two subsequent
calibration and attachment chapters identify and use these local inputs;
their witness identities do not alone prove the retained inequalities.

Fix a grade $d$ and its coefficient frame $V_d$. Let $\mathcal I_d$ be
the Maxwell equation ideal in the finite label space, and let $R_d$ be
the fixed grade-preserving campaign map of
Section~\ref{sec:marked-campaign-producer}. Denote its chosen local
section by $\mathcal V_d$, and the physical local section by
$\mathcal L_d$. For a coefficient $v\in V_d$, let
$\mathcal V_{N,p}(v;o)$ denote its normalized local insertion at $o$;
explicitly,
\[
 \mathcal V_{N,p}(v;o)=a_N^{-d}\,
        \mathcal V_d.\mathrm{places}(S_{\theta_N^{-1}}v;o),
\]
where $S_t$ is the factor-degree diagonal scaling and $\theta_N$ is the
original top threshold. Define $\mathcal L_{N,p}$ in the same way.
Both maps are complex-linear
in $v$. Their finite stencils, continuity, translation covariance and
signed symmetry are those of the fixed sections. Let
$X_{N,p}:V_d\to V_d$ be the terminal calibration, so that the physical
insertion of a reference label $x$ is
\[
                  O_{N,p}(x;o)=\mathcal L_{N,p}(X_{N,p}^{-1}x;o).
\]
All cutoffs below belong to the same selected family.

\subsection{The common calibration and its local error}

Apply the terminal invitation and calibration of
Section~\ref{sec:actual-marked-portfolio-calibration} to every frame
basis vector $e_j$, with the single fixed section $\mathcal V_d$.
The stage and delay purchases of
Section~\ref{sec:ex-marked-portfolio-calibration} supply the derivative
credit, its running-coupling upper bound, and the sphere inequality.
The selected-history attachment supplies the same terminal activity
keeper, transfer doors and finite-census estimates. Thus, eventually
in $N$ and for every $p$, there are columns $u_{j,N,p}$ with
\[
 |e_j-u_{j,N,p}|_1\le\delta_d,
                      \qquad \delta_d=(2\mathcal O_d^2)^{-1},
\]
where $\mathcal O_d$ is the fixed orbit cost. The same columns satisfy,
for every earlier opened pair of levels $(m,k)$ with
$\operatorname{ref}(d)\le k<N$ and $m\le N$, and for every signed
motion $g$,
\begin{equation}\label{eq:equation-column-root}
 \mathcal R_{m,N,p}\!\left(
  g\cdot\{\mathcal V_{N,p}(e_j;0)
                 -\mathcal L_{N,p}(u_{j,N,p};0)\}\right)
       \le c_{d,m,k}E_N e^{-s_N/8}.
\end{equation}
Here $s_N$ is the logarithmic scale, $s_N\to\infty$,
$\mathcal R_{m,N,p}$ is the centered root test, and $c_{d,m,k}$ is
the fixed opened-scale price. The amount $E_N$ is the explicit common
price in Section~\ref{sec:ex-marked-portfolio-calibration}. For every
$\eta>0$ it obeys $E_N\le e^{\eta s_N}$ eventually along any ordered
extraction. This uses the same subpower bounds on the credit, column
count, support exponent, insertion-site count, and the two reserve
factors stated there.

Average these columns over signed motions and conjugation, with the
twisted coefficient convention of the frame. Write the resulting map
as
\begin{equation}\label{eq:equation-common-calibration}
 P_{N,p}v=\sum_{(g,j)}\chi_{g,j}(v)\,g u_{j,N,p},
 \qquad
 v=\sum_{(g,j)}\chi_{g,j}(v)\,g e_j.
\end{equation}
The coefficients $\chi_{g,j}$ are complex-linear in $v$: for a conjugating
motion the coefficient of the inverse-transformed vector is conjugated
again. The map $P_{N,p}$ is therefore complex-linear. Reindexing the
finite average shows that it commutes with signed motions and conjugation.
Lemma~\ref{lem:ex-marked-calibration-inverse} gives
\begin{equation}\label{eq:equation-inverse-bound}
                       |P_{N,p}^{-1}v|_1\le2|v|_1.
\end{equation}
Use arbitrary invertible maps at the finitely many exceptional mesh
cutoffs, since all assertions using the columns are eventual.

For later use define a coefficient family $v_{N,p}$ to be subpower if,
for every $\eta>0$, its norm is eventually at most $e^{\eta s_N}$ on
the ordered filter under consideration. Every fixed linear coordinate
of such a family is subpower. Finite sums and products of subpower
nonnegative bounds are again subpower: split a target exponent between
the finitely many factors, and absorb fixed constants by a further
positive exponent margin. Equivalent norms on a fixed finite-dimensional
space give the same definition.

\subsection{A common tiling for the finite column sum}

We explain why \eqref{eq:equation-column-root} can be summed with
cutoff-dependent subpower coefficients. The bounds at opened pairs do
not have to hold at one common finest cell for all requests at once.
For any desired lower depth $w$, choose $m\ge\max(w,\operatorname{ref}(d))$
and then its purchased opening level $k\ge m$. They are fixed before
the cutoff tends to infinity. The price $c_{d,m,k}$ is consequently
constant, and the subpower bound on $E_N$ absorbs it. Since
$1/8-1/16=1/16>0$, \eqref{eq:equation-column-root} is eventually at most
$e^{-s_N/16}$. Fixed-stencil locality also gives chamber interiority
eventually. This yields centered comparisons at arbitrarily fine fixed
depths with that same decay rate.

To combine different column errors, transfer each such centered estimate
to every required larger cell. Fix a cell depth $m$ and a positive integer
margin denominator $J$. Choose a finer centered depth $k\ge m$ such that
$3J a_k\le a_m$, and at which the centered estimate holds. If its lattice
span is $h=a_k/a_N$, then a depth-$m$ cell has span $H=a_m/a_N$ and
$3h\le H/J$. The scales are aligned, so $h$ divides $H$.

At a point of the finer center grid, the full orbit of copies in the
larger tiling embeds injectively into cells of the finer centered tiling.
Its copied orientations are coordinate reflections, with the prescribed
complex conjugation. Translation invariance and reflection covariance
identify each of their centered root tests with the test at the origin.
Applying the planted inequality to the entire copy orbit and taking its
root bounds the larger-cell root by the same number. A bound on just
the unreplicated insertion would not suffice for this step.

For an arbitrary point in the larger cell with margin $3h$, bracket each
coordinate by consecutive fine-grid centers. The sixteen resulting
corners and the intervening coordinate segments lie in the margin $2h$.
The finer centered support therefore stays inside the larger cell along
all these segments. The reflection argument in
Section~\ref{ex-centered-estimates} gives the integer maximum principle
$r_j^2\le r_{j-1}r_{j+1}$ for the nonnegative root tests along one segment.
Its proof includes zero root tests. Applying it successively in four
coordinates bounds every intermediate root by the largest corner root.
Consequently every requested tiling of that larger side, and every
literal translated insertion at margin $H/J$, has root at most
$e^{-s_N/16}$. This argument applies to either allowed seam type.

Now let $Q_\alpha$ be the finitely many oriented column differences in
\eqref{eq:equation-common-calibration}, and let $r_\alpha(N,p)$ be
subpower scalar coefficients, constant in the spatial translation
variable. On any one of these larger tilings the root triangle inequality
gives
\[
 \mathcal R\left(\sum_\alpha r_\alpha Q_\alpha\right)
        \le\sum_\alpha|r_\alpha|\mathcal R(Q_\alpha)
        \le\left(\sum_\alpha|r_\alpha|\right)e^{-s_N/16}.
\]
For completeness, this triangle inequality follows by expanding the
product of copies in the root test of the sum. The planted inequality
bounds each term by the product of the individual root tests; summing
the expansion gives the corresponding power of their sum. Taking the
positive root proves the inequality, with the conjugated coefficient
copies contributing the same moduli. The finite sum of the coefficient
moduli is subpower, so the last display is at most $e^{-s_N/32}$
eventually. The comparisons have first been transferred to the same
tiling; intersecting unrelated sets of admissible cell depths would not
justify this addition.

Take $r_{g,j}=\chi_{g,j}(v_{N,p})$ for a subpower family $v$.
Linearity of the two sections and \eqref{eq:equation-common-calibration}
identify this sum, up to the harmless common sign, with
\begin{equation}\label{eq:equation-physical-difference}
       O_{N,p}(X_{N,p}P_{N,p}v_{N,p};o)
                              -\mathcal V_{N,p}(v_{N,p};o).
\end{equation}
It therefore has a centered discrepancy tending to zero, uniformly on
every fixed requested cell side and relative interior margin in the
ordered sense. The common stencil is the finite union of the two
section stencils and their finitely many orientations.

\subsection{Exact vanishing against arbitrary separated spectators}

The microscopic identity of Section~\ref{sec:marked-campaign-producer}
says that for every $z\in\mathcal I_d$, the local function read from
$R_dz$ has zero Wilson moment against every bounded measurable spectator
depending only on links outside a sufficiently large ball. Its radius
and minimum period can be chosen independently of $z$ at fixed $d$.
Indeed choose a basis of the finite-dimensional space $\mathcal I_d$,
take the maxima of its finitely many radii and minimum periods, and
expand $z$ in that basis. Continuity and the moment identity pass to
finite linear combinations. This also covers the zero ideal, using
the zero function.

Fix any finite spectator list, a distinguished slot, and a physical
separation $\delta>0$. The grades of the spectators are fixed, while
their coefficient values can be arbitrary. For each of those grades
the local physical section has a fixed stencil. Let $r$ bound their
lattice radii and let $s$ be the uniform microscopic radius just
obtained. Both are finite. For sufficiently large $N$, uniformly in
the period,
\[
                          a_N(s+r)<\delta,
\]
and the lattice side exceeds all the finite nonwrapping floors. If the
distinguished point is physically separated from each spectator point
by at least $\delta$, every link in every spectator stencil lies outside
the radius-$s$ ball about it. This follows from the triangle inequality:
the distance of a spectator link from the distinguished point is at
least the distance of its anchor minus $r$, which is greater than $s$.

The product of the other insertions is continuous on the compact
finite-volume configuration space, hence measurable and bounded. It is
local on those outside links. The microscopic identity therefore applies
to that entire product at once. Restoring the scalar normalization
$a_N^{-d}$ of the distinguished insertion preserves the value zero.
This proves exact vanishing at late finite cutoffs for every finite
spectator word. The cutoff threshold may depend on the grades and on
$\delta$, but not on the coefficients of the equation vector. The
coupling is nonzero on the selected Wilson history, as required by
the microscopic integration identity.

\subsection{From a root discrepancy to a spectator estimate}
\label{subsec:equation-spectator-price}

A vanishing root discrepancy still has to be converted to a bound on
ordinary correlations with spectators. Let $D_{N,p}(o)$ have the
common-tiling root price $y_{N,p}\to0$ obtained above. Fix a list of
$n\ge1$ insertion grades and coefficients, one distinguished slot,
and separation $\delta>0$. Choose an available cell depth $k$ with
$a_k<\delta$, and put $J=8(n+1)$. The centered observable bounds on
the same family give one finite budget $A\ge1$ for all the spectator
tokens and their finitely many reflected and conjugated versions,
at some finer depth $j\ge k$ with $3J a_j\le a_k$.

At late cutoffs let $L=a_k/a_N$ and $h=a_j/a_N$ be their integer
spans. Thus $3h\le L/J$. Shift the link tiling of side $L$ so that
every insertion lies at distance at least $\lfloor L/J\rfloor$ from
its chamber boundaries. Such a shift exists separately in each axis:
each of the $n$ coordinates excludes at most
$2\lfloor L/J\rfloor$ residues modulo $L$, and
$2n\lfloor L/J\rfloor<L$. Rebase the tiling so that the distinguished
point lies in cell zero. All insertions occupy distinct cells, since
two points in one cell have physical separation at most $a_k<\delta$.

For each spectator, reflect its position into the home cell and
transform its token by the inverse cell orientation, including ordinary
conjugation when prescribed. Its copy in the original cell is exactly
the original spectator. The centered budget at span $h$, the divisibility
of spans, and the margin $3h$ give interiority and root bound $A$ for
this home-cell function by the same planting and integer maximum
principle used above. The discrepancy insertion already has root bound
$y_{N,p}$ and interiority at its position in cell zero. Plant these
$n$ functions into their distinct cells to obtain
\begin{equation}\label{eq:equation-spectator-bound}
 \left|\mathbb E_{N,p}\left[
 D_{N,p}(o_i)\prod_{j\ne i}O_{N,p}(q_j;o_j)\right]\right|
                           \le A^{n-1}y_{N,p}.
\end{equation}
The normalized Wilson measure and its partition function are the same
for every factor. This estimate holds uniformly over the separated
tuples at late ordered cutoffs. The mesh threshold may depend on the
period. The constant $A$ depends on the fixed spectator list and the
chosen depths, which is sufficient for every individual test.

\subsection{The moving equation space and its limit}

The embedding of the campaign labels into the reference coefficient
space is literally
\[
                  E_{N,p}=X_{N,p}S_{\theta_N}\operatorname{coords}R_d.
\]
Define
\[
 B_{N,p}=X_{N,p}P_{N,p}X_{N,p}^{-1}E_{N,p},
                  \qquad K_{N,p}=B_{N,p}\mathcal I_d.
\]
All maps in this formula are the maps already selected above. The
calibrated terminal history and the symmetrized map $P_{N,p}$ imply
that $B_{N,p}$ intertwines conjugation and signed symmetries eventually.
Since the equation ideal has those symmetries, so does $K_{N,p}$.
The uniformly bounded equivariant projections of
Section~\ref{sec:finite-rank-kernel-transport} can therefore be chosen
with precisely this kernel. At the exceptional initial cutoffs use
the unsymmetrized bounded projections; only eventual symmetry is needed
for their limits.

Consider any row $x_{N,p}\in K_{N,p}$ whose inverse terminal coordinates
$X_{N,p}^{-1}x_{N,p}$ are subpower. Set
\[
                 v_{N,p}=P_{N,p}^{-1}X_{N,p}^{-1}x_{N,p}.
\]
The bound \eqref{eq:equation-inverse-bound} makes $v$ subpower. Also
$X_{N,p}v_{N,p}$ belongs to the original embedded campaign equation
space, directly from the definition of $K_{N,p}$. Equivalently there
is $z\in\mathcal I_d$ with
$v_{N,p}=S_{\theta_N}\operatorname{coords}R_dz$. In the normalized
section, $S_{\theta_N^{-1}}S_{\theta_N}=I$ cancels the threshold
exactly. Its value is consequently $a_N^{-d}$ times the microscopic
section for $R_dz$. Thus
\[
                         Z_{N,p}(o)=\mathcal V_{N,p}(v_{N,p};o)
\]
is an exactly vanishing comparison insertion of the preceding subsection.
It is measurable and bounded at every finite cutoff, and has a fixed
finite stencil. Equation~\eqref{eq:equation-physical-difference} becomes
$O_{N,p}(x_{N,p};o)-Z_{N,p}(o)$, whose discrepancy tends to zero.
This proves the local equation input for every such row in the moving
kernel, not just for a chosen finite basis of that kernel.

Finally take the common ordered extraction for the kernel projections
and correlations. Bounded rows in the limiting projection kernel can
be approximated by bounded rows of these finite kernels. Their inverse
terminal coordinates are subpower by Section~\ref{ex-column-subpower}.
For a compactly supported smooth spectator test on a separated domain,
its support has a positive separation margin. The corresponding
finite-cutoff sum with $Z_{N,p}$ is exactly zero eventually. The
spectator estimate \eqref{eq:equation-spectator-bound} and the
step-integral estimate of Section~\ref{ex-column-subpower} send the
difference from the physical insertion to zero. Passing to the ordered limit proves
that the limiting kernel annihilates every spectator test. This is the
all-test kernel inclusion used in
Section~\ref{sec:finite-rank-kernel-transport}. The separate finite rank
measurements then supply the opposite dimension bound.
